    \paragraph{Zhrnutie dosiahnutých výsledkov}:
    
    V tejto bakalárskej práci bola riešená úloha detekcie častí historických rukou písaných nomenklátorových šifrovacích kľúčov pomocou moderných metód umelej inteligencie. Hlavnou časťou práce bolo vytvorenie a spracovanie datasetu, ktorý zahŕňal prefotené historické dokumenty s štruktúrou, ktorá obsahuje všetky komponenty nomenklátoru. Pre detekciu boli využité modely architektúry YOLO11. Boli testované viaceré prístupy k spracovaniu dát ako binarizácia a augmentácia.

    Experimenty potvrdili, že vhodné predspracovanie dát zvyšuje presnosť tlačeného písmma. Binarizácia odstránila rušivé prvky a zvýraznila hrany. Augmentácia rozšírila obzor učenia a prispela k lepšej generalizácii modelu na odlišné dokumenty. 
    
    Najlepšie výsledky boli dosiahnuté pomocou YOLO11s.
    \paragraph{Prínos práce}:
    
    Aplikovanie moderných architektúr na detekciu špecifických časti nomenklátorových šifrovacích kľúčov, ktoré sú dôležitým zdrojom pre výskum dejín diplomacie a kryptológie. Pripravený dataset a navrhnutá metóda predspracovania môžu slúžiť ako základ pre ďalšie projekty v oblasti počítačového učenia zameraného na historické rukopisy.

    Ďalším prínosom je overenie účinnosti binarizácie a augmentácie pri práci s  rukou písanými dokumentami. Výsledky ukazujú, že aj pri relatívne malom množstve dát je možné dosiahnuť správnu detekciu, ak sú správne zvolené metódy prípravy a rozšírenia datasetu. Práca tak poskytuje cenné odporúčania pre ďalšie výskumné tímy a inštitúcie zaoberajúce sa analýzou historických dokumentov.
    \paragraph{Možnosti ďalšieho výskumu}:
    
    Na základe získaných výsledkov sa otvára viacero možnosti pre budúci výskum. Jednou z možností je rozšírenie datasetu o ďalšie typy historických šifrovacích dokumentov z rôznych geografických oblastí a časových období, čím by sa zvýšila robustnosť a univerzálnosť trénovaných modelov. Zároveň je vhodné experimentovať s pokročilejšími architektúrami, napríklad s modelmi zameranými na detekciu alebo s odličnými prístupmi, ktoré by mohli ešte lepšie, na prvý pohľad, nie veľmi odlišné vizuálne vzory historických rukopisov. Práca tak vytvára pevný základ hodný pre ďalšie inovácie.